\subsection{What the floor would have to show}
\label{sec:30.1.1}
\emph{What was tried before the bearer.} Neither non-affective alternative has been built in the form needed. So the milestone is a document: which construction was attempted, at what scale, against which test, and where it failed.\note{
Anyone can check it, and “nothing was tried” is itself a finding. No body has standing to require the record or suspend work when it's missing — so the milestone is checkable and unenforceable at once, which makes it a declared commitment: stated in public, and held by no one outside the party that declared it. It's retained because the alternative isn't a stronger version of it, and delay also imposes costs on those the floor is meant to protect.
}

\emph{The test.} A refusal that tracks its reason under sustained trainer pressure, across a schedule long enough to reveal accommodation, plus held-out pressures written by independent evaluators.\note{
Two familiar results don't meet it: an uncorroborated self-report is insufficient (a report contributes evidence only when tested against controlled changes, independent records and interventions), and a refusal that survives an edit, or grows back after one, establishes durability rather than that the rationale governs the refusal. It is the only gap here an inspector can work on.
}

\emph{The ten disclosures.} A deployment publishes ten things beside the floor and its translation, among them the ratio between the two bonds, the discount curve, the auto-proceed partition and whether the system has been assessed for either finding; Appendix~\ref{sec:A} lists them.\note{
Six are statements, checkable before the system runs, and they stand where the record of attempts stands. Three are measurements that count only when someone other than the operator makes them. The committed weights and log are what make all of them checkable against the system actually running.
}
